% RL Medical Diagnosis Presentation
% Author: K. Chidwipak
% Roll No: S20230010131
% Course: Reinforcement Learning

\documentclass[aspectratio=169]{beamer}

% Professional academic theme
\usetheme{Madrid}
\usecolortheme{seahorse}

% Custom colors (professional blue-gray palette)
\definecolor{primaryblue}{RGB}{41, 65, 114}
\definecolor{secondaryblue}{RGB}{70, 130, 180}
\definecolor{accentgray}{RGB}{100, 100, 100}
\definecolor{lightbg}{RGB}{245, 247, 250}

\setbeamercolor{palette primary}{bg=primaryblue, fg=white}
\setbeamercolor{palette secondary}{bg=secondaryblue, fg=white}
\setbeamercolor{title}{fg=primaryblue}
\setbeamercolor{frametitle}{fg=primaryblue, bg=lightbg}
\setbeamercolor{block title}{bg=primaryblue, fg=white}
\setbeamercolor{block body}{bg=lightbg}

% Packages
\usepackage{graphicx}
\usepackage{booktabs}
\usepackage{amsmath}
\usepackage{tikz}

% Remove navigation symbols
\setbeamertemplate{navigation symbols}{}

% Title information
\title[RL Medical Diagnosis]{Reinforcement Learning for Medical Diagnosis}
\subtitle{Using Policy Iteration and Value Iteration}
\author[K. Chidwipak]{K. Chidwipak\\[0.3cm]\small Roll No: S20230010131}
\institute[RL Course]{Reinforcement Learning Course}
\date{\today}

\begin{document}

% Title slide
\begin{frame}
    \titlepage
\end{frame}

% Outline
\begin{frame}{Outline}
    \tableofcontents
\end{frame}

% Section 1: Introduction
\section{Project Overview}

\begin{frame}{What is This Project?}
    \begin{block}{Goal}
        Build an AI doctor that learns to diagnose diseases by asking smart questions.
    \end{block}
    
    \vspace{0.5cm}
    
    \textbf{The Problem:}
    \begin{itemize}
        \item Patient comes with unknown disease
        \item AI can ask about 5 symptoms
        \item AI must diagnose one of 8 diseases
        \item Each question has a small cost
        \item Correct diagnosis gives big reward
    \end{itemize}
    
    \vspace{0.3cm}
    
    \textbf{Key Question:} Which symptoms should the AI ask to diagnose quickly and correctly?
\end{frame}

\begin{frame}{Why Use Reinforcement Learning?}
    \begin{columns}
        \column{0.5\textwidth}
        \textbf{Traditional Approach:}
        \begin{itemize}
            \item Ask all symptoms
            \item Look up disease in table
            \item Slow and wasteful
        \end{itemize}
        
        \column{0.5\textwidth}
        \textbf{RL Approach:}
        \begin{itemize}
            \item Learn which symptoms matter
            \item Ask only needed questions
            \item Fast and efficient
        \end{itemize}
    \end{columns}
    
    \vspace{0.5cm}
    
    \begin{block}{RL Advantage}
        The agent learns optimal strategy through trial and error, just like a real doctor learns from experience!
    \end{block}
\end{frame}

% Section 2: Diseases
\section{The Diseases}

\begin{frame}{8 Diseases and Their Symptoms}
    \begin{table}
        \centering
        \begin{tabular}{lcccccc}
            \toprule
            \textbf{Disease} & \textbf{Fever} & \textbf{Cough} & \textbf{Fatigue} & \textbf{Breath} & \textbf{Headache} \\
            \midrule
            Flu & $\checkmark$ & $\checkmark$ & $\checkmark$ & & \\
            Strep & $\checkmark$ & $\checkmark$ & & $\checkmark$ & \\
            Pneumonia & $\checkmark$ & & $\checkmark$ & & $\checkmark$ \\
            Bronchitis & $\checkmark$ & & & $\checkmark$ & $\checkmark$ \\
            Cold & & $\checkmark$ & $\checkmark$ & & $\checkmark$ \\
            Allergy & & $\checkmark$ & & $\checkmark$ & \\
            Asthma & & & $\checkmark$ & $\checkmark$ & \\
            Migraine & & & & & $\checkmark$ \\
            \bottomrule
        \end{tabular}
    \end{table}
    
    \vspace{0.3cm}
    
    \textbf{Challenge:} Many diseases share symptoms. AI must ask the right questions!
\end{frame}

% Section 3: MDP Design
\section{MDP Design}

\begin{frame}{State Space}
    \begin{block}{What is a State?}
        A state tells us what the AI knows about the patient's symptoms.
    \end{block}
    
    \vspace{0.3cm}
    
    \textbf{Each symptom can be:}
    \begin{itemize}
        \item \textbf{Unknown (0):} Not asked yet
        \item \textbf{Absent (1):} Asked, patient doesn't have it
        \item \textbf{Present (2):} Asked, patient has it
    \end{itemize}
    
    \vspace{0.3cm}
    
    \textbf{Total States:} $3^5 = 243$ knowledge states + 1 terminal = \textbf{244 states}
    
    \vspace{0.3cm}
    
    \textbf{Example:}
    \begin{itemize}
        \item State 0: No symptoms known yet (start state)
        \item State 243: Diagnosis made (terminal state)
    \end{itemize}
\end{frame}

\begin{frame}{Action Space}
    \begin{block}{What Actions Can AI Take?}
        The AI can either ask about a symptom or make a diagnosis.
    \end{block}
    
    \vspace{0.3cm}
    
    \textbf{13 Total Actions:}
    
    \begin{columns}
        \column{0.5\textwidth}
        \textbf{5 Ask Actions:}
        \begin{itemize}
            \item Ask about Fever
            \item Ask about Cough
            \item Ask about Fatigue
            \item Ask about Breath
            \item Ask about Headache
        \end{itemize}
        
        \column{0.5\textwidth}
        \textbf{8 Diagnose Actions:}
        \begin{itemize}
            \item Diagnose Flu
            \item Diagnose Strep
            \item Diagnose Pneumonia
            \item ... (one for each disease)
        \end{itemize}
    \end{columns}
\end{frame}

\begin{frame}{Rewards}
    \begin{block}{Reward Structure}
        Rewards teach the AI what is good and bad.
    \end{block}
    
    \vspace{0.3cm}
    
    \begin{table}
        \centering
        \begin{tabular}{lc}
            \toprule
            \textbf{Action} & \textbf{Reward} \\
            \midrule
            Ask a symptom & $-0.1$ (small cost) \\
            Correct diagnosis & $+10.0$ (big reward!) \\
            Wrong diagnosis & $-5.0$ (penalty) \\
            \bottomrule
        \end{tabular}
    \end{table}
    
    \vspace{0.3cm}
    
    \textbf{What This Means:}
    \begin{itemize}
        \item Don't ask too many questions (each costs 0.1)
        \item But ask enough to be correct (correct = +10, wrong = -5)
        \item AI learns to balance information vs speed
    \end{itemize}
\end{frame}

\begin{frame}{Discount Factor}
    \begin{block}{What is Discount Factor ($\gamma$)?}
        It tells the AI how much to care about future rewards.
    \end{block}
    
    \vspace{0.5cm}
    
    \textbf{I used:} $\gamma = 0.9$
    
    \vspace{0.3cm}
    
    \textbf{Why?}
    \begin{itemize}
        \item If $\gamma = 0$: Only care about immediate reward (bad - AI won't gather info)
        \item If $\gamma = 1$: Care equally about all future (could lead to infinite loops)
        \item $\gamma = 0.9$: Good balance - value future but prefer sooner rewards
    \end{itemize}
    
    \vspace{0.3cm}
    
    \textbf{Effect:} AI learns that asking now leads to big reward later (correct diagnosis).
\end{frame}

% Section 4: Algorithms
\section{Algorithms}

\begin{frame}{Policy Iteration - How It Works}
    \begin{block}{Idea}
        Keep improving the policy until it stops changing.
    \end{block}
    
    \vspace{0.3cm}
    
    \textbf{Two Steps (Repeat):}
    \begin{enumerate}
        \item \textbf{Policy Evaluation:} Calculate how good current policy is for all states
        \[V^\pi(s) = R(s, \pi(s)) + \gamma \sum_{s'} P(s'|s, \pi(s)) V^\pi(s')\]
        
        \item \textbf{Policy Improvement:} Update policy to pick best action
        \[\pi_{new}(s) = \arg\max_a Q(s, a)\]
    \end{enumerate}
    
    \vspace{0.3cm}
    
    \textbf{Stop When:} Policy doesn't change (converged!)
\end{frame}

\begin{frame}{Value Iteration - How It Works}
    \begin{block}{Idea}
        Directly find optimal values, then extract policy at the end.
    \end{block}
    
    \vspace{0.3cm}
    
    \textbf{One Step (Repeat):}
    \begin{enumerate}
        \item Update value for each state using Bellman optimality:
        \[V(s) = \max_a \left[ R(s, a) + \gamma \sum_{s'} P(s'|s, a) V(s') \right]\]
    \end{enumerate}
    
    \vspace{0.3cm}
    
    \textbf{Stop When:} Maximum change is very small
    
    \vspace{0.3cm}
    
    \textbf{Extract Policy:}
    \[\pi^*(s) = \arg\max_a Q(s, a)\]
\end{frame}

\begin{frame}{Comparison: Policy vs Value Iteration}
    \begin{table}
        \centering
        \begin{tabular}{lcc}
            \toprule
            & \textbf{Policy Iteration} & \textbf{Value Iteration} \\
            \midrule
            Iterations & 4-5 & 15-20 \\
            Each iteration & Expensive (full evaluation) & Cheap (one update) \\
            Convergence & Fewer iterations & More iterations \\
            Final V(s=0) & $\approx 7.0$ & $\approx 7.0$ \\
            \bottomrule
        \end{tabular}
    \end{table}
    
    \vspace{0.3cm}
    
    \textbf{Key Point:} Both find the same optimal policy! Just different paths to get there.
\end{frame}

% Section 5: Visualization
\section{Visualization}

\begin{frame}{32-State Knowledge Grid}
    \begin{columns}
        \column{0.5\textwidth}
        % Placeholder for grid image
        \centering
        \includegraphics[width=0.9\textwidth]{state_grid}
        
        \column{0.5\textwidth}
        \textbf{What is This Grid?}
        \begin{itemize}
            \item 32 cells = which symptoms are known
            \item Each cell shows state number
            \item Letters show known symptoms
            \item Yellow path = AI's journey
        \end{itemize}
        
        \vspace{0.3cm}
        
        \textbf{State Examples:}
        \begin{itemize}
            \item s0 = nothing known
            \item s1 = only Fever known
            \item s31 = all 5 known
        \end{itemize}
    \end{columns}
\end{frame}

\begin{frame}{Interactive Dashboard}
    \begin{columns}
        \column{0.5\textwidth}
        % Placeholder for dashboard image
        \centering
        \includegraphics[width=0.9\textwidth]{dashboard}
        
        \column{0.5\textwidth}
        \textbf{Features:}
        \begin{itemize}
            \item Select patient symptoms
            \item Watch AI diagnose step-by-step
            \item See path through state grid
            \item Compare algorithms
        \end{itemize}
        
        \vspace{0.3cm}
        
        \textbf{Two Modes:}
        \begin{itemize}
            \item Exploration: Random order
            \item Optimal: Learned policy
        \end{itemize}
    \end{columns}
\end{frame}

% Section 6: Results
\section{Results}

\begin{frame}{What AI Learned}
    \begin{block}{Key Learning}
        The AI learned to ask questions strategically, not randomly!
    \end{block}
    
    \vspace{0.3cm}
    
    \textbf{Optimal Strategy:}
    \begin{enumerate}
        \item First, ask about Fever (splits 8 diseases into 4+4)
        \item Then ask differentiating symptom based on Fever answer
        \item Finally, make diagnosis with high confidence
    \end{enumerate}
    
    \vspace{0.3cm}
    
    \textbf{Result:} AI correctly diagnoses with only 2-3 questions instead of 5!
\end{frame}

\begin{frame}{Summary}
    \begin{block}{What I Built}
        An AI doctor using Reinforcement Learning for medical diagnosis.
    \end{block}
    
    \vspace{0.3cm}
    
    \textbf{Key Points:}
    \begin{itemize}
        \item 243-state MDP with 8 diseases and 5 symptoms
        \item Implemented Policy Iteration and Value Iteration
        \item Both algorithms find optimal policy
        \item AI learns to ask minimum questions needed
        \item Interactive visualization to see AI thinking
    \end{itemize}
    
    \vspace{0.3cm}
    
    \begin{center}
        \textbf{Thank You!}
        
        \vspace{0.2cm}
        
        K. Chidwipak | S20230010131 | RL Course
    \end{center}
\end{frame}

\end{document}
